\documentclass[11pt]{article}
\usepackage[a4paper,margin=1in]{geometry}
\usepackage{amsmath,amssymb,mathtools}
\usepackage{booktabs,array}
\usepackage{enumitem}
\usepackage[dvipsnames]{xcolor}
\usepackage{hyperref}
\usepackage{microtype}
\usepackage{comment}

\hypersetup{colorlinks=true,linkcolor=MidnightBlue,urlcolor=MidnightBlue,
  pdftitle={Recitation 7: Bias, Standard Error, and Confidence Intervals},
  pdfauthor={Xuanxi Zhang}}

% Default: student handout. Define \WITHSOLUTIONS before inputting this file
% to include full answers and the instructor's 75-minute pacing guide.
\newif\ifsolutions
\ifdefined\WITHSOLUTIONS
  \solutionstrue
\else
  \solutionsfalse
\fi

\newcommand{\Pbb}{\mathbb{P}}
\newcommand{\E}{\mathbb{E}}
\newcommand{\Var}{\operatorname{Var}}
\newcommand{\Bias}{\operatorname{Bias}}
\newcommand{\MSE}{\operatorname{MSE}}
\newcommand{\se}{\operatorname{se}}
\newcommand{\rv}[1]{\widetilde{#1}}

\ifsolutions
  \newenvironment{answer}{%
    \par\medskip\noindent\textcolor{MidnightBlue}{\textbf{Solution.}}\ }
    {\par\medskip}
  \newcommand{\workingspace}[1]{}
\else
  \excludecomment{answer}
  \newcommand{\workingspace}[1]{\vspace{#1}}
\fi

\newcommand{\keybox}[1]{%
  \begin{center}
  \fcolorbox{MidnightBlue}{MidnightBlue!5}{%
    \begin{minipage}{0.91\textwidth}#1\end{minipage}}
  \end{center}}
\setlist[enumerate]{label=(\alph*),leftmargin=2em,itemsep=0.4em}
\setlist[itemize]{leftmargin=1.5em,itemsep=0.25em}

\title{\huge Recitation 7\\[0.35em]
  \normalsize Bias, Standard Error, and Confidence Intervals\\[0.35em]
  \large\ifsolutions Solutions and Teaching Notes\else Review and Practice Problems\fi}
\author{Xuanxi Zhang}
\date{October 23, 2026}

\begin{document}
\maketitle
\noindent\textit{Reading:}
\href{https://github.com/cfgranda/ps4ds/blob/main/ps4ds_preprint.pdf}
{\emph{Probability and Statistics for Data Science}},
Sections 9.2--9.3 and 9.8; Theorem 9.8, Exercise 9.3, and Example 9.45.

\section{Review: bias and standard error}
A parameter $\theta$ is fixed. Before observing the data, an estimator
$\rv{g}=h(\rv{x}_1,\ldots,\rv{x}_n)$ is random; after observing
$X=\{x_1,\ldots,x_n\}$, its value $g=h(x_1,\ldots,x_n)$ is a number.

\paragraph{Two different sources of error.}
For an estimator with a finite second moment,
\[
  \boxed{\Bias(\rv{g})=\E[\rv{g}]-\theta,\qquad
    \se(\rv{g})=\sqrt{\Var(\rv{g})}.}
\]
Bias is the average signed error; standard error measures variation across
repeated samples. Unbiased means $\E[\rv{g}]=\theta$, not that every
estimate equals $\theta$. Their contributions to mean squared error are
\[
  \MSE(\rv{g})=\E[(\rv{g}-\theta)^2]
    =\Bias(\rv{g})^2+\se(\rv{g})^2.
\]

\paragraph{The sample mean.}
If $\rv{x}_1,\ldots,\rv{x}_n$ are i.i.d. with mean $\mu$ and finite
variance $\sigma^2$, then
\[
  \rv{m}=\frac1n\sum_{i=1}^n\rv{x}_i,\qquad
  \E[\rv{m}]=\mu,\qquad
  \Var(\rv{m})=\frac{\sigma^2}{n},\qquad
  \se(\rv{m})=\frac{\sigma}{\sqrt n}.
\]
Linearity gives the mean; independence makes the cross-covariances zero
when computing the variance. Multiplying $n$ by four halves the SE.

\paragraph{Estimating the standard error from data.}
For $n\geq2$, define
\[
  m(X)=\frac1n\sum_{i=1}^n x_i,\qquad
  v(X)=\frac1{n-1}\sum_{i=1}^n(x_i-m(X))^2.
\]
The sample standard deviation $\sqrt{v(X)}$ describes the spread of the
\emph{observations}. The estimated SE of their mean is
$\se_{\mathrm{est}}=\sqrt{v(X)/n}$.
Problem 1 explains the denominator $n-1$.

\newpage
\section{Review: confidence intervals}
A $1-\alpha$ confidence procedure produces a random interval
$[\rv{\ell},\rv{u}]$ with coverage
\[
  \Pbb_\theta\{\rv{\ell}\leq\theta\leq\rv{u}\}\geq1-\alpha
\]
for each parameter value in the model. Exact equal-tail constructions may
have equality; normal approximations aim for coverage near the stated level.

\paragraph{The normal-approximation recipe.}
If $\rv{g}$ is approximately Gaussian and centered at $\theta$, then
\[
  \Pbb\!\left(\left|\frac{\rv{g}-\theta}{\se(\rv{g})}\right|
    \leq c_\alpha\right)\approx1-\alpha,
  \qquad c_\alpha=F_{\rv{z}}^{-1}(1-\alpha/2),
\]
where $\rv{z}\sim\mathcal N(0,1)$. Rearranging the event and, when
needed, estimating the SE gives the observed interval
\[
  \boxed{\bigl[g-c_\alpha\se_{\mathrm{est}},\,
                     g+c_\alpha\se_{\mathrm{est}}\bigr].}
\]
At 95\% confidence, use $c_{0.05}\approx1.96$.
The \textbf{margin of error} is the half-width $1.96\se_{\mathrm{est}}$.
For an i.i.d. sample with finite positive variance, the CLT motivates
\[
  \left[m(X)-1.96\sqrt{\frac{v(X)}n},\,
        m(X)+1.96\sqrt{\frac{v(X)}n}\right]
\]
when the sample is large enough for the normal approximation to be useful.

\paragraph{A population proportion.}
For i.i.d. $\rv{b}_i\sim\operatorname{Bernoulli}(\theta)$, let
$\rv{p}=n^{-1}\sum_i\rv{b}_i$. Then
\[
  \E[\rv{p}]=\theta,\qquad
  \se(\rv{p})=\sqrt{\frac{\theta(1-\theta)}n}.
\]
From an observed proportion $p$, a common estimate is
$\se_{\mathrm{est}}=\sqrt{p(1-p)/n}$. For planning before any data,
\[
  \theta(1-\theta)=\frac14-\left(\theta-\frac12\right)^2\leq\frac14
  \quad\Longrightarrow\quad
  \text{95\% margin of error}\ \leq\frac{0.98}{\sqrt n}
\]
under the normal approximation. This worst-case choice can be conservative.
The approximation needs enough expected successes and failures; it can be
poor near $\theta=0$ or $1$ with small samples.

\keybox{\textbf{Interpretation.}
Across repeated random samples, about 95\% of intervals from a calibrated
95\% procedure contain the fixed parameter. Once the interval is observed,
both its endpoints and the parameter are fixed; 95\% describes the
procedure, not a posterior probability for the parameter.}

\noindent\textit{Quick checks:} A larger sample reduces SE but does not
automatically remove sampling bias. Higher confidence gives a wider interval;
quadrupling the sample size approximately halves its width.

\newpage
\section*{Practice problems}
\noindent Three core problems; Problem 4 is a backup for a fast class.
The indicated times are for independent work before discussion.

\subsection*{Problem 1: Why divide by $n-1$? (Theorem 9.8)}
\noindent\textit{Independent work: 8 minutes.}
Let $a_1,\ldots,a_N$ be fixed population values with mean
$\mu_{\mathrm{pop}}=N^{-1}\sum_{k=1}^N a_k$ and variance
$\sigma_{\mathrm{pop}}^2=N^{-1}\sum_{k=1}^N(a_k-\mu_{\mathrm{pop}})^2$.
Draw $n\geq2$ samples independently and uniformly with replacement, and set
\[
  \rv{m}=\frac1n\sum_{i=1}^n\rv{x}_i,\qquad
  \rv{v}=\frac1{n-1}\sum_{i=1}^n(\rv{x}_i-\rv{m})^2.
\]
Prove that $\E[\rv{v}]=\sigma_{\mathrm{pop}}^2$.
\begin{enumerate}
  \item Show the identity
  \[
    \sum_{i=1}^n(\rv{x}_i-\rv{m})^2
    =\sum_{i=1}^n(\rv{x}_i-\mu_{\mathrm{pop}})^2
       -n(\rv{m}-\mu_{\mathrm{pop}})^2.
  \]
  \item Take expectations, using the mean and variance of $\rv{m}$.
    Where is independence needed?
  \item What is the bias if we divide the squared-deviation sum by $n$
    instead of $n-1$?
\end{enumerate}
\workingspace{7.5cm}

\begin{answer}
\textbf{(a)} Expand
$\rv{x}_i-\rv{m}=(\rv{x}_i-\mu_{\mathrm{pop}})
-(\rv{m}-\mu_{\mathrm{pop}})$. Since
$\sum_i(\rv{x}_i-\mu_{\mathrm{pop}})=n(\rv{m}-\mu_{\mathrm{pop}})$,
the cross term is $-2n(\rv{m}-\mu_{\mathrm{pop}})^2$; adding the
$n$ squared mean terms gives the identity.

\textbf{(b)} Uniform sampling gives
$\E[\rv{x}_i]=\mu_{\mathrm{pop}}$ and
$\Var(\rv{x}_i)=\sigma_{\mathrm{pop}}^2$. Independence yields
$\Var(\rv{m})=n^{-2}\sum_i\Var(\rv{x}_i)=\sigma_{\mathrm{pop}}^2/n$.
Because $\E[\rv{m}]=\mu_{\mathrm{pop}}$,
\[
  \E\!\left[\sum_{i=1}^n(\rv{x}_i-\rv{m})^2\right]
  =n\sigma_{\mathrm{pop}}^2
    -n\E[(\rv{m}-\mu_{\mathrm{pop}})^2]
  =(n-1)\sigma_{\mathrm{pop}}^2.
\]
Dividing by $n-1$ proves $\boxed{\E[\rv{v}]=\sigma_{\mathrm{pop}}^2}$.
The same proof works for any i.i.d. sample with finite variance.

\textbf{(c)} Let $\rv{q}=n^{-1}\sum_i(\rv{x}_i-\rv{m})^2$.
Then
\[
  \E[\rv{q}]=\frac{n-1}{n}\sigma_{\mathrm{pop}}^2,\qquad
  \boxed{\Bias(\rv{q})=-\frac{\sigma_{\mathrm{pop}}^2}{n}.}
\]
Centering at the fitted sample mean reduces the expected squared-deviation
sum by one population variance, which explains the correction.
\end{answer}

\newpage
\subsection*{Problem 2: Participation (Exercise 9.3)}
\noindent\textit{Independent work: 10 minutes.}
A poll records party preference and stated intention to vote:
\begin{center}
\begin{tabular}{@{}lrr@{}}
\toprule
Intention & Democrat & Republican\\
\midrule
Certain to vote & 20 & 30\\
Likely to vote & 40 & 10\\
Not likely to vote & 30 & 20\\
\bottomrule
\end{tabular}
\end{center}
A separate study gives turnout probabilities $0.8$, $0.1$, and $0$ for
these three intention categories, respectively.
\begin{enumerate}
  \item Let $\alpha_R$ be the population fraction in the joint category
    ``certain to vote and Republican.'' Assume independent, uniform
    sampling with replacement. Construct an unbiased estimator,
    derive its true standard error, and compute its observed value.
    Also give a numerical estimate of the SE.
  \item Estimate the Republican share \emph{among those who actually
    vote}, using the turnout probabilities. State the assumptions that
    make this adjustment appropriate.
\end{enumerate}
\workingspace{9cm}

\begin{answer}
\textbf{(a)} Define $\rv{b}_i=1$ if respondent $i$ is in the stated
joint category, and $0$ otherwise. Then the $\rv{b}_i$ are i.i.d.
Bernoulli$(\alpha_R)$, so the estimator
$\rv{a}_R=150^{-1}\sum_{i=1}^{150}\rv{b}_i$ satisfies
\[
  \E[\rv{a}_R]=\alpha_R,\qquad
  \boxed{\se(\rv{a}_R)=\sqrt{\frac{\alpha_R(1-\alpha_R)}{150}}.}
\]
The observed estimate is $a_R=30/150=\boxed{0.20}$. Substituting this
value for the unknown parameter gives
\[
  \se_{\mathrm{est}}=\sqrt{\frac{0.20(0.80)}{150}}
     \approx\boxed{0.03266}\quad\text{(3.27 percentage points).}
\]
This estimates the true SE; it is not its known exact value.

\textbf{(b)} The weighted counts are
\[
  \text{Republican: }0.8(30)+0.1(10)=25,\qquad
  \text{Democrat: }0.8(20)+0.1(40)=20.
\]
Normalize by the total weighted turnout:
\[
  \boxed{\text{Estimated Republican vote share}
    =\frac{25}{25+20}=\frac59\approx55.56\%.}
\]
Assume representative sampling, accurate and stable party preferences,
and turnout rates that apply to this population and are the same for both
parties within each intention category. Under these assumptions, the
ratio estimates $\Pbb(\text{Republican}\mid\text{votes})$.
The raw preference share is $60/150=40\%$, which targets a different
quantity. A ratio of unbiased component estimates need not be unbiased.
\end{answer}

\newpage
\subsection*{Problem 3: Planning the sample size (Example 9.45)}
\noindent\textit{Independent work: 8 minutes.}
We estimate a population prevalence $\theta$ by testing $n$ people
selected independently and uniformly with replacement. Treat test results
as accurate i.i.d. Bernoulli$(\theta)$ measurements. We want a 95\%
confidence interval with margin of error at most $0.01$
(\emph{one percentage point}). Use the normal approximation with $1.96$.
\begin{enumerate}
  \item Without knowing $\theta$, use $\theta(1-\theta)\leq1/4$ to
    determine a sufficient sample size for this planning rule.
  \item If we plan at $\theta=0.05$, how many tests does the rule require?
    Why is this smaller, and does it work for every possible prevalence?
  \item For the sample size in (b) and true $\theta=0.05$, let
    $\rv{t}$ count the positive tests. Identify its distribution and the
    integer values for which $|\rv{t}/n-0.05|\leq0.01$.
    Write the exact probability of meeting this error tolerance as a
    finite sum; no numerical evaluation by hand is required.
\end{enumerate}
\workingspace{10cm}

\begin{answer}
\textbf{(a)} The worst-case margin is $1.96/(2\sqrt n)=0.98/\sqrt n$.
Thus
\[
  \frac{0.98}{\sqrt n}\leq0.01
  \quad\Longleftrightarrow\quad n\geq98^2,
  \qquad \boxed{n=9604}.
\]
This is a sufficient \emph{normal-approximation} planning value, not an
exact finite-sample coverage guarantee.

\textbf{(b)} At $\theta=0.05$, the variance is $0.05(0.95)=0.0475$,
which is much smaller than $1/4$. Therefore
\[
  n\geq\frac{1.96^2(0.05)(0.95)}{0.01^2}=1824.76,
  \qquad \boxed{n=1825}.
\]
Always round up. This calculation assumes the planning prevalence is
correct; it does not ensure the desired margin at other prevalences
with larger variance.

\textbf{(c)} Here $\rv{t}\sim\operatorname{Binomial}(1825,0.05)$, and
\[
  0.04\leq\frac{\rv{t}}{1825}\leq0.06
  \quad\Longleftrightarrow\quad73\leq\rv{t}\leq109.5
  \quad\Longleftrightarrow\quad73\leq\rv{t}\leq109.
\]
The exact probability is
\[
  \boxed{\sum_{t=73}^{109}\binom{1825}{t}(0.05)^t(0.95)^{1825-t}
    \approx0.95319.}
\]
So the probability of an error of at most one percentage point is about
95.32\%, close to the nominal 95\%.
\end{answer}

\newpage
\subsection*{Problem 4 (backup): A polling margin}
\noindent\textit{Independent work: 5 minutes.}
In a random sample of $1500$ voters, $55\%$ support a proposition and
$45\%$ oppose it. Model the responses as independent Bernoulli variables
with support probability $\theta$. The population support-minus-opposition
margin is $d=\theta-(1-\theta)=2\theta-1$.
\begin{enumerate}
  \item Express the margin estimator $\rv{d}$ in terms of the sample
    support proportion $\rv{p}$. Is it unbiased for $d$?
  \item Find its true standard error and estimate that SE from the data.
  \item Construct an approximate 95\% confidence interval for $d$.
    Explain why treating the support and opposition proportions as
    independent would give the wrong standard error.
\end{enumerate}
\workingspace{11cm}

\begin{answer}
\textbf{(a)} Since the opposition proportion is $1-\rv{p}$,
\[
  \rv{d}=\rv{p}-(1-\rv{p})=2\rv{p}-1,\qquad
  \E[\rv{d}]=2\theta-1=d.
\]
The observed margin is $2(0.55)-1=0.10$, or 10 percentage points.

\textbf{(b)} Variance scales by the square of the multiplier, so
\[
  \Var(\rv{d})=4\Var(\rv{p})
    =\frac{4\theta(1-\theta)}{1500},\qquad
  \se(\rv{d})=2\sqrt{\frac{\theta(1-\theta)}{1500}}.
\]
The estimated SE is
\[
  \se_{\mathrm{est}}=2\sqrt{\frac{0.55(0.45)}{1500}}
    \approx0.02569\quad\text{(2.57 percentage points).}
\]

\textbf{(c)} The approximate interval is
\[
  0.10\pm1.96(0.02569)
    \approx\boxed{[0.0496,\,0.1504]}.
\]
This corresponds to a support-minus-opposition margin between about
4.96 and 15.04 percentage points. Both proportions use the same respondents:
the opposition proportion is determined by the support proportion.
Their covariance is $-\Var(\rv{p})$, so ignoring it underestimates
the variance of their difference.
\end{answer}

\vfill
\noindent\textit{Source:} adapted from the Fall 2025 recitation source
\texttt{mine/r3/r3.tex}, Section 7.9. The unbiasedness question has been added.

\ifsolutions
\newpage
\section*{Instructor pacing guide (75 minutes)}
\begin{center}
\begin{tabular}{@{}lll@{}}
\toprule
Time & Activity & Suggested use\\
\midrule
0--8 min & Bias and standard error & Estimator vs. observed estimate; MSE\\
8--15 min & Mean and proportion & Derive SE; distinguish SD from SE\\
15--25 min & Confidence intervals & Normal recipe; coverage; half-width\\
25--42 min & Theorem 9.8 & 8 min individual work, 9 min proof\\
42--59 min & Exercise 9.3 & 10 min individual work, 7 min discussion\\
59--73 min & Example 9.45 & 8 min individual work, 6 min discussion\\
73--75 min & Wrap-up & Three short checks below\\
\bottomrule
\end{tabular}
\end{center}
Use Problem 4 only if the class finishes early, with five minutes of
independent work before discussion. Keep the three core work periods.

\paragraph{Board plan for the variance proof.}
Start from the identity in Problem 1(a), then take expectations term by
term. Emphasize that $\rv{x}_i$ and $\rv{m}$ are \emph{not} independent;
independence is used between different sample observations when computing
$\Var(\rv{m})$. The residuals sum to zero, leaving only $n-1$ free
residuals. This gives intuition for the denominator; the expectation
calculation is the proof. Unbiasedness of $\rv{v}$ does not imply
unbiasedness of $\sqrt{\rv{v}}$.

\paragraph{Teaching points for Exercise 9.3.}
Part (a) targets a joint intention/preference category among all people,
so use denominator 150. The true SE contains the unknown $\alpha_R$;
the numerical SE uses the observed fraction. Part (b) changes the target
to a conditional proportion among voters, so normalize weighted counts
by 45. The turnout assumption is conditional independence of turnout and
party preference given the intention category. The weighted counts are
predictions, not observed election totals.

\paragraph{Teaching points for Example 9.45.}
The error tolerance $0.01$ is an absolute proportion error, not a 1\%
relative error. Distinguish a variance upper bound from a planning guess,
and always round sample sizes upward. The exact binomial check is for
$\theta=0.05$ only; evaluating the sum can be supplied by the instructor.
The preprint's integer upper limit 110 should be 109 for the stated
tolerance, and its displayed ``0.959\%'' has an incorrect percent sign.
The corrected event used here has probability $0.95319$.

\paragraph{Closing checks.}
Why does the sample variance divide by $n-1$? How does quadrupling $n$
change SE and CI width? What is random in a confidence interval statement?

\paragraph{Sources and adaptations.}
The review follows Carlos Fernandez-Granda's
\emph{Probability and Statistics for Data Science}, Sections 9.2--9.3
and 9.8. Problem 1 gives a centered-sum proof of Theorem 9.8
(pp.\ 370--371); Problem 2 rephrases Exercise 9.3 (p.\ 431), preserving
the data and adding a numerical SE estimate; Problem 3 turns Example 9.45
(pp.\ 410--411) into guided practice and corrects its binomial check.
Problem 4 reuses the polling-margin exercise from the user's Fall 2025
source, \texttt{mine/r3/r3.tex}, with an added bias calculation.
\fi

\end{document}
